\section{刑事和解（公诉）}
\begin{enumerate}
  \item 核心。
        \begin{enumerate}
          \item \bs{当事人互让互谅 \yjt 法院负责组织、主持}。\al{负责调解 \cuo}。
          \item \bs{和解 \bdy 无罪。各自在职权范围内从宽处理}。
        \end{enumerate}
  \item 条件。\bs{民间、人财、3 年下，7 年过失渎职外}。民间纠纷；人身、民主权利，财产权利；判 3 年下。
        \begin{enumerate}
          % \item 你认错，他原谅，都愿意。侵害特定人，事实清。
          \item \bs{5 年不犯故意罪}。
          \item 不属于民间纠纷。\bs{雇、黑、寻、多次聚众斗殴和伤害}。雇凶伤人、黑社会性质组织犯罪、寻衅滋事、多次聚众斗殴/故意伤害。
          \item 各阶段都可以提。
        \end{enumerate}
  \item 公安。侦查阶段和解，\bs{不撤案}，向检提从宽处理建议。
  \item 检察院。\yi 审查起诉中。向法院建议从宽处罚 \hz 可以酌定不起诉。\er 审理中。没他啥事。
  \item 法院。审理中。\yi 可以从宽处罚，但\zd{不能宣告无罪}。\er 部分和解，部分适用。
  \item 被害人，嫌疑人/被告人。2 方主体。
        \begin{enumerate}
          \item 被害人\bs{死了叫直接和解，活着叫代为和解}（人在才有的代）。
          \item 人民调解员、辩护人、诉代、当事人亲友均可参加促成和解。
          \item 民事可协商，专属公检法的不能。\al{甲可就对乙免除刑事处分达成和解 \cuo，也可就民事赔偿达成和解 \dui}。
        \end{enumerate}
  \item 和解协议。
        \begin{enumerate}
          \item \zd{不盖公章}只有双方签名、没有强制力、不能申请强制执行、只能申请确认效力。
          \item 原则\zd{一次履行，即时履行}。例外。提供担保。
          \item \zd{履行完毕产生终结力，再诉不支持}。
        \end{enumerate}
\end{enumerate}
